With the rapid advancement of global technology, the intelligence level of electronic devices has risen significantly alongside the development of integrated circuits. In daily life, as the core module of various electronic products, switching power supplies have garnered increasing societal attention. While chip development has spurred the emergence of diverse software, the progress of battery materials lags. Concurrently, the public’s demands for charger power and charging speed are escalating. Flyback AC/DC switching power supply converters, valued for their excellent isolation, simple design, and low cost, thrive in medium-to-low-power fast-charging domains. However, their large losses and low switching frequencies can no longer keep pace with electronic device advancements. To satisfy surging societal needs and explore more potential of flyback converters, this paper designs a flyback switching power supply converter chip based on an asymmetric half-bridge topology for medium-power fast-charging applications.
The flyback converter with a secondary-side feedback structure proposed herein optimizes existing issues in flyback converters. Multiple control modes are introduced to achieve high conversion efficiency and low power consumption across the full load range. Standby power consumption is reduced via turning off redundant modules and lowering switching frequency in burst mode. The newly proposed Zero-Voltage Resonant Valley Switching (ZV-RVS) mode incorporates soft-start time to introduce reverse magnetization current, enabling zero-voltage turn-on of the high-side power transistor and optimizing light-load efficiency. Under heavy loads, the critical conduction mode controls asymmetric complementary conduction of power transistors to achieve soft switching, meeting high-output-power demands. Additionally, a valley-locking circuit is designed to eliminate switching frequency fluctuations caused by valley-skipping within the power overlap range. A peak-current control circuit is added to ensure correspondence between secondary-side feedback voltage and load current under a wide output range, fulfilling multi-mode control requirements.
For further optimization of the Zero-Voltage Resonant Valley Switching mode, a precise valley-conduction technique is proposed. Based on traditional quasi-resonant technology, this technique adopts a more accurate sampling scheme. Through a step-by-step approximation strategy, it eliminates delay errors from the drive circuit, aiming to turn on the low-side power transistor at the lowest point of the resonant valley. This realizes minimal switching loss of the Zero-Voltage Resonant Valley Switching mode under light loads while addressing reliability issues related to inductor current spikes.
Concerning high-efficiency control of the asymmetric half-bridge flyback converter, a dynamic demagnetization-time calibration technique is proposed. An integral ramp circuit with dynamic compensation controls the conduction time of the low-side power transistor, gradually overlapping it with the energy-transfer time of the excitation inductor. This avoids increased conduction losses and reliability issues caused by excessively long or short conduction times of the low-side power transistor, further optimizing the circuit’s conversion efficiency.
The converter chip designed in this paper is developed using SMIC 0.18um BCD technology. Operating under an input range of 90–264Vac, it features a rated output range of 5–20V and an output load range of 0–5A. Simulation tests demonstrate that, under the multi-mode control strategy and various key technologies, the circuit significantly enhances system conversion efficiency at a maximum output power of 100W. The peak conversion efficiency reaches 90\%, and the load regulation rate is 1.12\%, meeting performance index requirements.